\documentclass[10pt,a4paper]{amsart}
\usepackage[margin=19mm]{geometry}
\usepackage{amsmath,amssymb,mathtools}
\usepackage[scr=rsfs,cal=euler]{mathalfa}
\usepackage{xfrac}
\usepackage[unicode,hidelinks,bookmarks=false]{hyperref}
\newcommand{\ie}{i.e.\ }
\newcommand{\cf}{cf.\ }
\newcommand{\resp}{respectively\ }
\newcommand{\Subsection}{\subsection}
\providecommand{\nobreakdash}{\nobreak\hbox{-}}
\setlength{\emergencystretch}{2em}
\multlinegap0pt
\usepackage{bm}
\usepackage{bbm}
\usepackage{dsfont}
\usepackage{xspace}
\usepackage{cancel}
\usepackage{mathtools}
\usepackage{amsthm}
\makeatletter
\@ifundefined{theorem}{\newtheorem{theorem}{Theorem}}{}
\@ifundefined{corollary}{\newtheorem{corollary}[theorem]{Corollary}}{}
\@ifundefined{lemma}{\newtheorem{lemma}[theorem]{Lemma}}{}
\makeatother
\title[A theory of generalized Lam\'e curves]{A theory of generalized Lam\'e curves}
\author{You-Cheng Chou, Chin-Lung Wang, Po-Sheng Wu}
\date{Source: arXiv:2604.21880v2 (CC BY 4.0)}
\begin{document}
\maketitle

\noindent\emph{Source and licence.} A theory of generalized Lam\'e curves. Version arXiv:2604.21880v2 (CC BY 4.0), distributed under the Creative Commons Attribution 4.0 International licence, as recorded in the arXiv licence metadata for this exact version and shown on the abstract page https://arxiv.org/abs/2604.21880v2. Full rights notice: licenses/arxiv-2604.21880-CC-BY-4.0.txt.\par\smallskip

\noindent\textit{Editorial scope.} Counted unit: the Theorem whose source label is \texttt{thm\_deg} in the article (source0.tex lines 1523--1533, source label \texttt{thm\_deg}), delivered together with the article's own complete original proof of it (source0.tex lines 1534--1630, one \texttt{proof} environment of 97 lines). No mathematics is written, repaired or re-derived: statement and proof are the article's own text line for line. Required context retained uncounted: lemma\_boundary (lines 1076--1092); cor\_boundary (lines 1176--1186); special-pt (lines 1330--1337). Every definition, display and label of those lines is retained unchanged. Omitted from this package and not redistributed: 1--1075, 1093--1175, 1187--1329, 1338--1522, 1631--3446. Nothing inside a retained excerpt is omitted or altered: every retained body is delivered exactly as the article's lines are written, so no omission receipt of any kind accompanies this package. Citations: the retained excerpts carry no citation command at all, so no bibliography entry of the article is transcribed and no reference is redistributed. Reference adaptation: none was needed and none was made; every reference inside the delivered unit resolves inside it, so no unresolved reference or citation is produced. Printed slips kept literal: no printed word, symbol, hypothesis or number of the counted statement or of its proof was altered. Layout and class: the article's own class is \texttt{amsart} with packages including inputenc, amsmath, palatino, mathpazo, amsfonts, amssymb, mathtools, eucal, xy, datetime, amsthm, comment, multicol, graphicx; the wrapper is the lane's pinned \texttt{amsart} editorial wrapper at 10pt on A4, which restates the theorem-like environments the retained text needs and adds the article's own macro definitions that the retained text needs (none), with extra wrapper packages (bm, bbm, dsfont, xspace, cancel, mathtools). The delivered numbering is the unit's own. Assets: no retained span contains a figure environment, a graphics-inclusion command, a PDF-page inclusion, a table or a TikZ picture; no image file is copied into this package, the package declares no asset dependency and no image file is redistributed with it. Licence: version arXiv:2604.21880v2 of this article is distributed under the Creative Commons Attribution 4.0 International licence, as recorded in the arXiv licence metadata for this exact version and shown on the abstract page https://arxiv.org/abs/2604.21880v2. A full rights notice is delivered at licenses/arxiv-2604.21880-CC-BY-4.0.txt. Characters: every retained excerpt is pure ASCII, so the delivered engine, XeLaTeX with its default Latin Modern text font, reports no missing character.
\begin{lemma} \label{lemma_boundary}
The system of equations
\begin{equation} \label{equation_boundary}
    f_{\mu}(\alpha) := \sum_{\nu=1, \neq \mu}^k \frac{1}{\alpha_{\mu} - \alpha_{\nu}} - \frac{(k-1)x}{\alpha_{\mu}} + \sum_{\nu=1, \neq \mu}^k \frac{x}{\alpha_{\nu}} =0, \qquad \mbox{for $\mu =1 ,\dots,k$}
\end{equation}
has a unique solution (up to permutation and scaling) with $\alpha_\mu \neq \alpha_\nu$ for $x \in \mathbb{C}\setminus\{ 0, \frac{1}{2k}, \dots, \frac{k-2}{2k} \}$. 

Moreover, the Jacobian of $f$ evaluated at the corresponding solution, $J_f(\alpha)$,
\[
\begin{split}
    J_f(\alpha)_{\mu\mu} &= \sum_{\nu=1,\neq \mu}^k \frac{-1}{(\alpha_\mu-\alpha_\nu)^2} + \frac{(k-1)x}{\alpha_\mu^2}, 
    \\
    J_f(\alpha)_{\mu\nu} &= \frac{1}{(\alpha_\mu - \alpha_\nu)^2} - \frac{x}{\alpha_\nu^2} \text{ for } \mu \neq \nu,
\end{split}
\]
has rank $k-1$ for $x \in \mathbb{C}\setminus\{0, \frac{1}{2k}, \dots, \frac{k-1}{2k}\}$, and rank $k-2$ for $x = \frac{k-1}{2k}$. 
\end{lemma}

\begin{corollary} \label{cor_boundary}
Let $\alpha_1=1$ and $f$ be the system in Lemma~\ref{lemma_boundary}. There exist Zariski open sets 
\[
\begin{split}
    U &\subset \left\{ (x,\alpha)\in \mathbb{C}^{k+1} | \alpha_1=1; \ \alpha_i \neq \alpha_j   \right\},
    \\
    V &\subset \left\{ (x, \boldsymbol{\rho}) \in \mathbb{C}^{k+1} \Big| \sum_{\mu=1}^k \rho_\mu =0 \right\},
\end{split}
\]
such that $\Big( (\mathbb{C}\setminus\{ 0, \frac{1}{2k}, \dots,\frac{k-2}{2k} \}), \mathbf{0} \Big) \subset V$ and $f: U \rightarrow V$ is a morphism of degree $k!$.
\end{corollary}

\begin{theorem}[Classification of Type-I boundary] \label{special-pt}
Given $\mathbf{n} \in (\mathbb{C} \setminus \{0\} )^r$ with total weight $n\in \mathbb{N}$,
the Type-I boundary (where $h=\infty$) of $\overline{\mathcal{Y}}_{\mathbf{n},\mathbf{p}}(\tau)$ consists of the points $[(\mathbf{a}, \infty)]$ satisfying
\[
    \{[\mathbf{a}]\} = \{[p_1]^{k_1}, \dots, [p_r]^{k_r}\},
\]
where the multiplicities $k_j$ satisfy $k_j \leq 2n_j$ whenever $n_j \in \frac{1}{2}\mathbb{N}$.
\end{theorem}

\begin{theorem} \label{thm_deg}
Let $I := \{ i\ |\ n_i \in \tfrac{1}{2}\mathbb{N} \}$. The addition map
\[
\sigma_{\mathbf{n},\mathbf{p}}: \overline{Y}_{\mathbf{n},\mathbf{p}}(\tau) \to E_\tau
\] 
is generically of degree 
\[
    \sum_{J \subset I} (-1)^{|J|} \binom{n_J +r}{r+1}_{\rm comb},
\]
where $n_J := n - \sum_{j \in J}(2n_j +1)$, and $\binom{A}{B}_{\rm comb} := 0$ if $A<B$.
\end{theorem}

\begin{proof}
It is equivalent to study the generic degree of $\sigma_{\mathbf{n},\mathbf{p}}$ on the open set $Y_{\mathbf{n}, \mathbf{p}}(\tau)$, since the boundary has strictly lower dimension.
The proof proceeds in two steps. First, we compute the intersection number of the divisors $D_{\mu,\rho_\mu}$ in $E_\tau^n$ for general $\boldsymbol{\rho}$. By deforming to $|\rho_\mu| \to \infty$, we combinatorially count the points and obtain $n! \binom{n+r}{r+1}$. Second, we analyze the limit $\rho_\mu \to 0$ and apply an inclusion-exclusion argument to systematically subtract degenerate solutions (where $a_\mu \to p_i$), ensuring all remaining intersection points are valid. Because $Y_{\mathbf{n},\mathbf{p}}(\tau)$ is defined in the symmetric product ${\rm Sym}^n E_\tau$, the generic degree of $\sigma_{\mathbf{n},\mathbf{p}}$ is $1/n!$ times the number of valid solutions.

Given $\boldsymbol{\rho}=(\rho_1,\dots,\rho_n) \in \mathbb{C}^n$ with $\sum_{\mu=1}^n \rho_\mu=0$, we define 
\[
D_{\mu,\rho_\mu} := \left\{ \sum_{i = 1}^{r} \sum_{\nu = 1, \neq \mu}^n 
n_s\Big(
\zeta(a_\mu - a_\nu) - \zeta(a_\mu - p_i) + \zeta(a_\nu - p_i)
\Big) =  \rho_\mu \right\} \subset E_{\tau}^n,
\]
for $\mu=1,\dots,n$. Furthermore, given $c \in E_\tau$, we define
\[
Z^c(\boldsymbol{\rho}) :=  \left( \bigcap_{\mu=1}^n D_{\mu,\rho_\mu}\right) \cap \sigma_{\mathbf{n},\mathbf{p}}^{-1}(c) \subset E_\tau^n.
\]

We claim that, for generic $c$ and $\boldsymbol{\rho}$, by counting with multiplicities, 
\[
\#(Z^c(\boldsymbol{\rho})) = n! \binom{n+r}{r+1}.
\]
Assume $|\rho_\mu| >M>0$ for all $\mu$. For $M$ large enough, the points of $Z^c(\boldsymbol{\rho})$ cluster into components corresponding to partitions of the indices:
\[
Z^c(\boldsymbol{\rho}) = \bigsqcup_{\substack{I_0 \sqcup \dots \sqcup I_r = \{1,\dots,n\} \\ I_0 \neq \emptyset }} P_{I_0,\dots,I_r}^c (\boldsymbol{\rho}).
\]
As $M \to \infty$, each component $P^c_{I_0,\dots,I_r}$ is characterized by its limit:
\[
P_{I_0,\dots,I_r}^c(\infty) = 
\left\{
\mathbf{a}\in E_\tau^n \;\middle|\; 
\begin{array}{l}
a_\mu = a_\nu \text{ for all } \mu, \nu \in I_0, \\
a_\mu = p_j \text{ for all } \mu \in I_j, \forall j \in \{1,\dots,r\}
\end{array}
\right\}
\]
with multiplicity
\[
(|I_0|-1)! \prod_{i=1}^r |I_i|!.
\] 
The multiplicity corresponds to the number of ways to obtain $P_{I_0,\dots,I_r}^c(\infty)$ from the intersection of the irreducible components of $D_{\mu,\rho_\mu}$ for $\mu=1,\dots,n$. We also note that $P_{I_0,\dots,I_r}^c(\infty)$ is the disjoint union of $|I_0|^2$ points, corresponding to the number of $|I_0|$-torsion points in $E_\tau$.

Combining all the information above, we have 
\[
\begin{split}
\#(Z^c(\boldsymbol{\rho})) &=  \sum_{I_0 \sqcup \dots \sqcup I_r = \{1,\dots,n\}} | P_{I_0,\dots,I_r}^c (\boldsymbol{\rho}) | \cdot (|I_0|-1)! \prod_{i=1}^r |I_i|!
\\
&= \sum_{I_0 \sqcup \dots \sqcup I_r = \{1,\dots,n\}}   |I_0|\prod_{i=0}^r |I_i|!
\\
&= \sum_{|I_0|+\dots+|I_r|=n} \frac{n!}{\prod_{i=0}^r |I_i|!} |I_0|\prod_{i=0}^r |I_i|!
\\
&= n! \sum_{|I_0|+\dots+|I_r|=n} |I_0| = n! \binom{n+r}{r+1}.
\end{split}
\]
This finishes the first step.

Next, let $\boldsymbol{\rho} \to 0$. We study the intersection points that degenerate, i.e., where $a_\mu \to p_i$. We claim that the set of degenerate solutions can be stratified by subsets $J \subset I := \{ i \mid n_i \in  \tfrac{1}{2} \mathbb{N} \}$ satisfying 
\[
n_J := n - \sum_{j \in J} (2n_j+1) \ge 0.
\]
More precisely, for $J=\{ 1,\dots,l \}$, up to permutation of coordinates, the degenerate solutions in $E_\tau^n$ take the form:
\[
    \mathbf{a} = (\underbrace{p_1, \dots, p_1}_{2n_1+1}, \dots, \underbrace{p_l, \dots, p_l}_{2n_l+1}, a'_1, \dots, a'_{n_J}),
\]
where the remaining coordinates form a point $\{[a'_1], \dots, [a'_{n_J}]\} \in Y_{\mathbf{n}',\mathbf{p}}^{c'}(\tau)$,
with $\mathbf{n}' = ( -n_1-1,\dots,-n_l-1,n_{l+1}, \dots,n_r )$ and 
\[
    Y_{\mathbf{n}',\mathbf{p}}^{c'}(\tau) := Y_{\mathbf{n}',\mathbf{p}}(\tau) \cap \sigma_{\mathbf{n}',\mathbf{p}}^{-1}(c') \text{ with }  c' = c- \sum_{j=1}^l (2n_j+1) p_j.
\]
Counting with multiplicity, there are exactly $n! \deg (\sigma_{\mathbf{n}',\mathbf{p}}: Y_{\mathbf{n}',\mathbf{p}}(\tau) \to E_\tau)$ such solutions in $E_\tau^n$.

To verify this, consider the special case $J=\{1\}$. As in the proof of Theorem~\ref{special-pt}, let $m$ be the number of coordinates $\mu$ for which $a_\mu \to p_1$. Without loss of generality, let 
\[
    a_\mu = p_1 + \sum_{k \ge 1} \alpha_{\mu k} \epsilon^k, \quad \text{for } \mu=1,\dots, m.
\]
Comparing the coefficients of $\epsilon^{-1}$ gives:
\[
\begin{cases}
    \displaystyle\sum_{\nu=1}^{m} \frac{1}{\alpha_{\nu 1}} = 0, \\[15pt]
    \displaystyle\sum_{\substack{\nu=1 \\ \nu \neq \mu}}^{m} \frac{1}{\alpha_{\mu 1} - \alpha_{\nu 1}} - \frac{n_1-1}{\alpha_{\mu 1}} + \sum_{\substack{\nu=1 \\ \nu \neq \mu}}^m \frac{1}{\alpha_{\nu 1}}  = 0.
\end{cases}
\]
By Lemma~\ref{lemma_boundary}, the solution set $\{ \alpha_{\mu 1} \}$ exists if and only if $m = 2n_1+1$. The constant terms for the remaining equations (for $\mu = 2n_1+2,\dots,n$) precisely impose the condition
\[
    \{[a_{2n_1+2}], \dots, [a_n]\} \in Y_{\mathbf{n}',\mathbf{p}}(\tau). 
\]
The higher-order terms $\{ \alpha_{\mu k} \}$ can be solved inductively by Corollary~\ref{cor_boundary}. Finally, the number of such degenerate solutions in $E_\tau^n$ is given by
\[
\begin{split}
    \binom{n}{2n_1+1} (2n_1+1)!  (n-2n_1-1)! & \deg ( \sigma_{\mathbf{n}',\mathbf{p}}: Y_{\mathbf{n}',\mathbf{p}}(\tau) \to E_\tau ) 
    \\
     = n! & \deg ( \sigma_{\mathbf{n}',\mathbf{p}}: Y_{\mathbf{n}',\mathbf{p}}(\tau) \to E_\tau ).
\end{split}
\]
Here, the binomial coefficient corresponds to the choice of the $2n_1+1$ coordinates converging to $p_1$, the $(2n_1+1)!$ factor corresponds to the number of valid leading-order solutions $\{ \alpha_{\mu 1} \}$, and the final term corresponds to the ordered degree of $\sigma_{\mathbf{n}',\mathbf{p}}$ on $E^{n-2n_1-1}_\tau$. 

The proof generalizes to an arbitrary subset $J = \{ 1,\dots,l \}$. To correctly isolate the non-degenerate solutions, we derive the degree formula in the inclusion-exclusion form. We complete the proof.
\end{proof}

\end{document}
